\chapter{Discusión}
\label{ch:discusion}

\section{Sostenibilidad Intertemporal y Reglas Fiscales}

Los resultados del Capítulo \ref{ch:resultados} exigen matices más finos que una confirmación o refutación binaria de las hipótesis del Capítulo \ref{ch:introduccion}. La evidencia no permite sostener, con significatividad convencional, que el Estado argentino incrementó sistemáticamente su superávit primario frente al aumento de la ratio Deuda/PIB durante 2004--2025 ($\rho=-0.0071$, $p=0.564$). Esta falta de significatividad estadística, unida al signo puntual negativo del parámetro, constituye evidencia empírica a favor de la Hipótesis 1: la política fiscal argentina no cumplió con la regla de reacción activa requerida por la condición de transversalidad de \citet{bohn1998}.

El diagnóstico de validez instrumental de la estimación IV-2SLS matiza, sin revertir, esta lectura. Los instrumentos utilizados (VIX y el diferencial de riesgo soberano regional basado en el ETF EMB) resultan estadísticamente sólidos en la primera etapa ($F=13.00$, por encima del umbral de 10 pero con menor margen que el hallado con la serie de contingencia de EMBI+), el test de Wu-Hausman confirma la endogeneidad del riesgo soberano y el test de sobreidentificación de Sargan no rechaza la validez conjunta de los instrumentos ($0.608$, $p=0.436$). Con un diseño instrumental válido, el coeficiente del EMBI+ instrumentado ($0.0017$) alcanza significatividad marginal al 10\% ($p=0.076$): una señal débil, pero ya no una ausencia de canal. Como advierten \citet{blanchard2021,blanchard2022}, en economías con historial de cesación de pagos y descalce de monedas resulta complejo aislar variables verdaderamente exógenas al ciclo fiscal doméstico, respondiendo la percepción de riesgo a factores idiosincráticos y eventos de renegociación; esa dificultad de identificación es compatible tanto con una señal débil genuina como con ruido de muestra pequeña, y esta tesis no fuerza una lectura por sobre la otra.

El signo negativo del coeficiente ($\rho=-0.0071$), aun sin significatividad estadística, indica que ante incrementos en la deuda heredada el resultado primario no se ajustó al alza. Esta inoperancia de la regla de Bohn traslada el ajuste hacia canales exógenos a la decisión fiscal discrecional (devaluación real, licuación inflacionaria y quitas forzosas), tal como documenta el Capítulo \ref{ch:descriptivo}. La Sección \ref{sec:fatiga_discusion} retoma esta lectura y la articula con el Nivel I del marco teórico.

Corresponde precisar el alcance temporal de este diagnóstico. La reestimación por subperíodos con especificación dinámica con ampliación de rezagos (DOLS; Sección \ref{sec:sensibilidad_temporal}) arroja un $\rho$ positivo en las cuatro especificaciones y significativo en tres de ellas (entre 0.036 y 0.060). El parámetro nulo de la muestra completa refleja una respuesta inestable entre ventanas temporales que no satisface la exigencia de sistematicidad de \citet{bohn1998}, en consonancia con los quiebres estructurales múltiples de la Sección \ref{sec:quiebres_estructurales} del Capítulo \ref{ch:descriptivo}.

El diagnóstico DOLS no queda aislado: el VECM que lo reemplaza como técnica de referencia sobre la ventana ampliada 1996--2025 (Sección \ref{sec:resultados_vecm} del Capítulo \ref{ch:resultados}), a pedido del director de tesis, converge en el mismo resultado nulo por una vía metodológicamente distinta. Tratar las cuatro variables del sistema como conjuntamente endógenas, en lugar de purgar la endogeneidad de corto plazo con adelantos y rezagos \textit{ad hoc} como hace DOLS, no revela una reacción fiscal estabilizadora oculta tras el sesgo de la técnica de ecuación única: la velocidad de ajuste del resultado primario hacia el equilibrio de largo plazo ($\alpha_{pb}=-0.00002$, $p=0.996$) permanece no significativa bajo la especificación de referencia, y en cuatro de cinco especificaciones alternativas de rezagos. Es la propia deuda, no el resultado primario, la que carga con la velocidad de ajuste significativa del sistema ($\alpha_{deuda}=-0.0736$, $p=0.004$), un resultado que anticipa la lectura de esta tesis sobre el canal dominante de ajuste, desarrollada más abajo. Que dos técnicas de estimación con supuestos y tratamientos de la endogeneidad tan distintos, DOLS de ecuación única con ampliación dinámica y VECM de sistema con corrección de error, coincidan en el mismo diagnóstico nulo, sobre dos ventanas muestrales que comparten apenas 88 de 120 observaciones, es la evidencia más robusta de esta tesis a favor de $H_{1a}$: no es un artefacto de una técnica particular.

Esa convergencia no es, sin embargo, inmune a la fragilidad documentada en el Capítulo \ref{ch:resultados}: el propio contraste de cointegración que sostiene el VECM en la ventana ampliada depende de excluir el año 1999 de la muestra, y el pedido explícito del director de extender la ventana ``aun con cambios estructurales extremos'' produjo, en ese sentido, precisamente el resultado que cabía esperar: una relación de largo plazo cuya existencia formal es más frágil, no más sólida, cuanto más se extiende la muestra hacia el tramo de mayor inestabilidad estructural de la serie.

\section{La Hipótesis de Fatiga Fiscal}
\label{sec:fatiga_discusion}

El modelo de umbral de \citet{hansen1999,hansen2000}, estimado sobre la serie real de EMBI+, identifica un punto de quiebre candidato en $\tau^*=449$ puntos básicos, con una reacción fiscal estimada menor en el régimen de estrés ($\beta_2 = 0.0126$, $p=0.385$) que en el régimen ``normal'' ($\beta_1 = 0.0305$, $p=0.171$). Este ordenamiento cualitativo coincide con la hipótesis de fatiga fiscal de \citet{ghosh2013}. A diferencia del diagnóstico que arrojaba la serie de contingencia de EMBI+ usada hasta la Sección \ref{sec:auditoria_fallback} del Capítulo \ref{ch:metodologia}, que no rechazaba linealidad, el contraste \textit{bootstrap} (Sup-LM, $p=0.076$) ahora sí la rechaza, al 10\% pero no al 5\%. Esta señal exige, sin embargo, una calificación inmediata: el umbral hallado cae cerca del extremo inferior del rango de búsqueda y particiona la muestra de manera extremadamente desigual (14 trimestres en el régimen ``normal'' contra 74 en ``estrés''), la condición en la que un contraste de umbral de series de tiempo resulta menos confiable, y explica por qué ninguno de los dos coeficientes de régimen alcanza significatividad individual pese al rechazo de linealidad en conjunto. La evidencia en niveles es, en síntesis, una señal débil y frágil, no una confirmación.

El modelo de \citet{ghosh2013} fue calibrado, además, para economías avanzadas con acceso continuo a financiamiento en moneda doméstica. Su aplicación a Argentina, donde la restricción vinculante es la disponibilidad de divisas y no la curva de Laffer fiscal interna (Nivel I del marco teórico, Capítulo \ref{ch:marco_teorico}), somete la hipótesis a una especificación potencialmente inadecuada para el contexto bimonetario \citep{eichengreen2003}.

La especificación en primeras diferencias ($\Delta pb_t$), detallada en la Sección \ref{sec:robustez_hansen_dpb}, ofrece una evidencia más contundente y menos frágil que la de niveles. Al utilizar la serie en diferencias, que satisface el supuesto de estacionariedad exigido por \citet{hansen2000}, el \textit{bootstrap} Sup-LM rechaza la linealidad con holgura ($p<0.001$, frente a $p=0.019$ que arrojaba la serie de contingencia) y ambos coeficientes de régimen resultan significativos con la ordenación esperada. A diferencia de la especificación en niveles, aquí la partición muestral (74 contra 14, en sentido inverso a la de niveles) no concentra el resultado en un puñado de observaciones del régimen que sostiene la significatividad de los coeficientes individuales. La evidencia de un umbral es, en consecuencia, sensible a la especificación elegida, pero ya no unilateralmente débil: la especificación de primera diferencia, la que satisface el supuesto formal del test, ofrece la evidencia más sólida de fatiga fiscal no lineal de todo el capítulo. Las restricciones del tamaño muestral ($n=88$) y la volatilidad del indicador de riesgo (desvío estándar de 1230 pbs sobre la serie real, frente a los 876 que exhibía la serie de contingencia) constituyen limitaciones reales del diseño que obligan a interpretar ambos resultados con cautela.

El diagnóstico de esta tesis se ancla en el Nivel I del marco teórico (Capítulo \ref{ch:marco_teorico}): los ajustes del stock de deuda operaron mediante canales exógenos a la decisión fiscal discrecional. Esta interpretación es coherente con la literatura macroeconómica estructuralista \citep{aruguete2021, damill2015}, según la cual en economías bimonetarias la sostenibilidad está acotada por la restricción externa de divisas.

Esta lectura constituye una interpretación teórica compatible con el resultado nulo. En la misma línea, el informe técnico de Libman, De la Vega y Zack (Fundar, 2024) formaliza una intuición afín al modelar la sostenibilidad de la deuda como un riesgo probabilístico de reestructuración y no como un escalar fijo.

La cuantificación directa de la contribución de cada canal al Ajuste Stock-Flujo ($SF_t$), mediante la descomposición contable de la identidad $d_t = \frac{1+r_t}{1+g_t}d_{t-1} - pb_t + SF_t$ (Sección \ref{sec:restriccion_intertemporal}) para los episodios de mayor variación discreta de la ratio Deuda/PIB (2005, 2018 y 2020; Capítulo \ref{ch:descriptivo}), permitiría transformar esta interpretación en evidencia directa. Como primera aproximación, la Tabla \ref{tab:sft_ilustrativo} aplica esa descomposición usando $d_t$, $d_{t-1}$ y $pb_t$ del dataset de la tesis, junto con una tasa de interés real efectiva $r_t$ aproximada a partir de series del FMI.

\begin{table}[H]
\centering
\caption{\label{tab:sft_ilustrativo} Descomposición Ilustrativa del Ajuste Stock-Flujo ($SF_t$), Episodios Seleccionados.}
\begin{tabular}{lrrrrrr}
\toprule
\textbf{Año} & \textbf{$d_t$} & \textbf{$d_{t-1}$} & \textbf{$pb_t$} & \textbf{$r_t$ (aprox.)} & \textbf{$g_t$} & \textbf{$SF_t$ implícito} \\
\midrule
2005 (reestructuración) & 70.5\% & 106.0\% & 5.7\% & $-7.3\%$ & 8.9\% & $-14.1\%$ \\
2018 (depreciación) & 90.3\% & 76.4\% & $-9.5\%$ & $-21.4\%$ & $-2.6\%$ & $19.2\%$ \\
2020 (quita + licuación) & 105.6\% & 93.7\% & $-10.4\%$ & $-27.6\%$ & $-9.9\%$ & $20.0\%$ \\
\bottomrule
\end{tabular}
\begin{flushleft}
\footnotesize \textit{Nota.} $d_t$, $d_{t-1}$, $pb_t$: dataset propio de la tesis. $r_t$: aproximación construida a partir de series del \textit{IMF WEO DataMapper} (gobierno general). $g_t$: crecimiento real del PIB. Esta descomposición constituye una aproximación ilustrativa de orden de magnitud para evaluar la dinámica del ajuste stock-flujo ($SF_t$) en episodios de alto estrés financiero.
\end{flushleft}
\end{table}

Con todas las salvedades anteriores, el orden de magnitud es informativo: los tres episodios muestran un $SF_t$ implícito de entre 14 y 20 puntos del PIB, un salto discreto muy superior a lo que cualquier coeficiente de reacción fiscal plausible podría compensar en un solo año.

\section{Transmisión Macrofiscal y Gestión de Pasivos}
\label{sec:implicancias_svar_cir}

La descomposición de varianza del SVAR Restringido (Sección \ref{sec:svar_resultados} del Capítulo \ref{ch:resultados}), estimada sobre la ventana ampliada de treinta años (1996--2025, $n=120$) que reemplaza a la ventana 2004--2025 como referencia del sistema, revierte la lectura que este capítulo sostenía sobre veintiséis años de datos. Allí, los shocks al resultado primario explicaban el $47.5\%$ de la varianza del error de pronóstico de la deuda a 20 trimestres y los shocks a la brecha del producto un $28.6\%$ adicional, frente a una contribución conjunta de apenas $11.5\%$ del tipo de cambio real y el riesgo soberano, y un $12.4\%$ de la inercia propia de la deuda. Sobre la ventana de treinta años, la inercia propia pasa a explicar el $54.2\%$ de la varianza de la deuda, el canal cambiario y de riesgo soberano conjunto ($25.8\%$) supera ahora al fiscal y de actividad conjunto ($20.0\%$), y ningún shock individual domina de forma aislada como lo hacía antes el superávit primario. Esta reversión, verificada como hallazgo genuino y no como artefacto de la extrapolación del TCRM de 1996 (Sección \ref{sec:svar_resultados} del Capítulo \ref{ch:resultados}), obliga a matizar, no a abandonar, el argumento de este capítulo. Sigue sin haber evidencia de que sea una regla de reacción fiscal estabilizadora la que mueve la deuda período a período: el propio shock al resultado primario cae de primer a segundo lugar en importancia relativa, y DOLS y el VECM ya documentan por vías independientes la ausencia de esa regla. Pero la explicación proximal deja de concentrarse en los desbalances primarios y el ciclo real, y pasa a repartirse entre la inercia de la propia deuda y el canal cambiario-riesgo soberano, un canal que remite de manera más directa que el fiscal al Paradigma de la Restricción Externa (Nivel I del marco teórico) que esta tesis privilegia como explicación sustantiva. En ese sentido, el hallazgo fortalece, más que debilita, el argumento central de que el ajuste de la deuda argentina transita por canales exógenos a la decisión fiscal discrecional: incluso la porción de la dinámica que antes se atribuía al resultado primario resulta, con treinta años de evidencia, menos determinante que la inercia propia del stock de deuda y que el riesgo soberano y cambiario. Esto tampoco contradice la evidencia sobre los canales del Ajuste Stock-Flujo (Sección \ref{sec:hoja_balance} del Capítulo \ref{ch:descriptivo}): una cosa son las perturbaciones que mueven la deuda período a período dentro del sistema SVAR, ahora sobre todo su propia inercia y el riesgo soberano-cambiario, y otra distinta son los canales por los que se resolvieron históricamente los episodios de sobreendeudamiento, a saber la depreciación real, la licuación inflacionaria y las reestructuraciones con quita, un punto que se retoma en las implicancias de economía política. Para el debate de política, la lectura es, si acaso, más cauta que antes: la consolidación fiscal primaria sigue siendo condición necesaria, pero opera sobre una dinámica en la que el ciclo real pesa ahora menos, y el riesgo soberano, el canal cambiario y la propia inercia de la deuda pesan más, de lo que sugería la ventana original.

\section{Implicancias de Economía Política}

Aun reconociendo estas limitaciones estadísticas, los coeficientes hallados son compatibles con los planteos de economía política sobre restricciones estructurales: la dificultad de sostener un esfuerzo fiscal primario suficiente remite a la tensión entre las exigencias financieras del Tesoro y las condiciones socioeconómicas de la población \citep{cantamutto2020}. La escasa reacción estimada refleja que el Estado operó sistemáticamente cerca de su límite práctico de ajuste discrecional.

\section{Proyecciones Probabilísticas post-2025 y Gráficos de Abanico}

El horizonte posterior a 2025 demanda abandonar los modelos contables deterministas, dado que la evidencia del Capítulo \ref{ch:resultados} no permite calibrar con precisión una regla de reacción fiscal para incorporar en la identidad de acumulación de deuda. Ante esta limitación, se optó por un Análisis de Sostenibilidad de la Deuda (DSA) que combina un componente determinista (tres sendas con supuestos macroeconómicos explícitos) y uno estocástico (Monte Carlo, 1.000 iteraciones), siguiendo el estándar de organismos multilaterales \citep{imf2003,imf2021}. La matriz de varianzas y covarianzas del componente estocástico deriva de la covarianza de residuos del SVAR restringido (Sección \ref{sec:dsa_resultados} del Capítulo \ref{ch:resultados}); una especificación anterior, calibrada a mano y luego reestimada con desvíos GARCH(1,1) y con la correlación condicional dinámica de un modelo DCC-GARCH \citep{engle2002}, se desarrolló y validó en una etapa previa de este trabajo y queda documentada aparte, en un memo de trabajo para el director de tesis, no en este documento.

A diferencia de las aplicaciones estandarizadas de DSA que asumen de forma irrestricta la simetría y normalidad de las perturbaciones macroeconómicas, esta tesis reconoce que las series temporales de la economía argentina exhiben propiedades estadísticas de colas pesadas y asimetría extrema, coherentes con crisis de balanza de pagos recurrentes: un patrón ya documentado en la Sección \ref{sec:propiedades_muestra} del Capítulo \ref{ch:descriptivo} para el EMBI+, y confirmado aquí para las propias perturbaciones del componente estocástico del DSA. Concretamente, se estimaron los grados de libertad de una distribución $t$ de Student multivariada por método de momentos ($\nu \approx 4.8$, a partir de la curtosis en exceso muestral: $\nu = 4 + 6/\text{curtosis}$). Se prefiere este estimador al ajuste por máxima verosimilitud conjunta de los tres parámetros de la $t$ (\textit{loc}, \textit{scale}, $\nu$): sobre la serie real de TCRM, ese ajuste conjunto da $\nu=1.54$ para el resultado primario y $\nu=2.61$ para el tipo de cambio real, ambos por debajo de 2, con varianza teóricamente infinita: una solución de esquina arrastrada por un puñado de trimestres genuinamente extremos (devaluación de 2016, crisis de 2018, ajuste de 2023--2024) y no por un error de datos, pero sí un patrón de inestabilidad conocido del MLE conjunto sin restringir en muestras chicas con observaciones extremas aisladas (\texttt{codigo/modelos/fase17\_calibracion\_nu\_dsa.py}). Los valores de $\nu\approx4.8$ obtenidos por método de momentos son, en cambio, acordes a la curtosis en exceso observada en las perturbaciones históricas de resultado primario ($25.3$) y de tipo de cambio real ($4.5$, sobre la serie real de TCRM; la serie de contingencia usada hasta la Sección \ref{sec:auditoria_fallback} del Capítulo \ref{ch:metodologia} arrojaba $14.4$). Este ajuste metodológico, que preserva la volatilidad calibrada de cada variable pero añade una probabilidad mayor de eventos extremos, dota a los gráficos de abanico resultantes de un realismo analítico superior frente a la subestimación del riesgo de insolvencia fiscal que caracteriza a los modelos simétricos tradicionales.

\subsection{Escenarios Deterministas}

Partiendo de una ratio inicial de 74\% del PIB en 2025, un valor calibrado a partir del nivel de deuda consolidada del SPNF de cierre observado en el Capítulo \ref{ch:descriptivo} (70--71\% en los últimos trimestres de 2025 según el panel trimestral), la Tabla \ref{tab:dsa_determinista} resume la trayectoria de la ratio Deuda/PIB bajo tres escenarios hacia 2035: Optimista (superávit primario de 2.5\% del PIB, crecimiento de 4.5\%), Referencia (superávit de 1.5\%, crecimiento de 3.5\%) y Estrés (contracción del producto, superávit exiguo de 0.5\% y una depreciación real acumulada del 25\%).

\begin{table}[H]
\centering
\caption{\label{tab:dsa_determinista} DSA determinista: Deuda Pública Consolidada Ampliada (SPNF + BCRA) / PIB, escenarios seleccionados (\%).}
\begin{tabular}{lrrr}
\toprule
\textbf{Año} & \textbf{Optimista} & \textbf{Referencia} & \textbf{Estrés} \\
\midrule
2026 & 74.0 & 74.0 & 74.0 \\
2030 & 61.8 & 77.6 & 246.3 \\
2033 & 52.9 & 80.7 & 571.9 \\
2034 & 50.0 & 81.7 & 753.8 \\
2035 & 47.1 & 82.8 & 992.1 \\
\bottomrule
\end{tabular}
\begin{flushleft}
\footnotesize \textit{Nota.} Elaboración propia, simulación determinista ($\alpha=0.463$ proporción de deuda en moneda extranjera). La fila 2035 corresponde a una iteración anual adicional bajo los supuestos inalterados de cada escenario, preservando la trazabilidad a la ventana prospectiva de diez años (2026--2035) declarada en el Capítulo \ref{ch:metodologia}.
\end{flushleft}
\end{table}

\begin{figure}[H]
\centering
\caption{Trayectorias del DSA Determinista bajo los Tres Escenarios (2026--2035)}
\label{fig:dsa_determinista}
\includegraphics[width=0.85\textwidth]{dsa_determinista.png}
\par\vspace{2mm}
{\footnotesize\textit{Nota.} Trayectorias correspondientes a la Tabla \ref{tab:dsa_determinista}. Elaboración propia.}
\end{figure}

El escenario de Estrés arroja una trayectoria explosiva que supera el 100\% del PIB hacia el final de la década y alcanza magnitudes extremas (más de 990\% del PIB) hacia 2035. Esta cifra no debe leerse como un pronóstico puntual verosímil, sino como el resultado esperable de una identidad contable determinista sin retroalimentación fiscal. Su valor analítico radica en ilustrar la sensibilidad de la trayectoria de deuda a la simultaneidad de perturbaciones adversas, el fenómeno que \citet{celasun2006} identifican como la falla central de los enfoques deterministas, y no en constituir una proyección literal del sendero más probable.

\subsection{Simulación Estocástica: Gráfico de Abanico de Monte Carlo}

La especificación de referencia deriva la matriz de covarianza del componente estocástico directamente de los residuos estimados del SVAR restringido sobre la ventana ampliada de treinta años (Sección \ref{sec:dsa_resultados} del Capítulo \ref{ch:resultados}). Una especificación anterior, calibrada a mano y luego reestimada con GARCH(1,1) univariado y con un modelo DCC-GARCH, se desarrolló y validó en una etapa previa de este trabajo, con una probabilidad de insolvencia hacia 2035 del orden de $31$--$33\%$; queda documentada aparte, en un memo de trabajo para el director de tesis, pendiente de que defina cuál de las dos especificaciones integrar a la versión final.

\begin{figure}[H]
\centering
\caption{Gráfico de Abanico: Distribución de Probabilidad de la Trayectoria de Deuda/PIB}
\label{fig:fan_chart_final}
\includegraphics[width=0.85\textwidth]{dsa_grafico_abanico.png}
\par\vspace{2mm}
{\footnotesize\textit{Nota.} Simulación de Monte Carlo (1.000 iteraciones, parámetros del escenario de Referencia con perturbaciones multivariadas correlacionadas, distribución $t$ de Student con $\nu\approx4.8$, matriz de covarianza derivada del SVAR restringido). Bandas de 10--90\% y 25--75\%; línea central: mediana. Elaboración propia.}
\end{figure}

A diferencia del escenario de Estrés determinista, la simulación estocástica reportada en la Figura \ref{fig:fan_chart_final} introduce perturbaciones anuales aleatorias alrededor de los parámetros centrales del escenario de Referencia, generando 1.000 trayectorias alternativas de la ratio Deuda/PIB hacia 2035. La mediana de estas trayectorias se ubica en $82{,}1\%$ del PIB hacia 2035 (percentil 10--90: $52{,}6\%$--$120{,}3\%$; percentil 25--75: $65{,}6\%$--$101{,}0\%$), mientras que las bandas de confianza se ensanchan progresivamente con el horizonte de proyección.

El resultado cuantitativo central es que la probabilidad de que la ratio Deuda/PIB supere la frontera crítica del 100\% del producto hacia 2035 es del $\mathbf{26{,}1\%}$: aproximadamente uno de cada cuatro escenarios simulados. Esta frontera del 100\% del PIB coincide con el pico histórico de la ratio de deuda argentina documentado en el Capítulo \ref{ch:descriptivo} (2020). Esta cifra debe interpretarse como una medida de la vulnerabilidad estructural pasiva de la economía argentina frente a la volatilidad histórica de sus fundamentos, y no como una probabilidad neta que ya incorpore una eventual respuesta fiscal correctiva; corresponde exclusivamente al escenario de Referencia, ejecutado sobre la matriz de covarianza fundada en el SVAR. A diferencia de la especificación anterior, no se dispone todavía de una réplica de este ejercicio centrada en los escenarios Optimista y de Estrés ni de un desglose de sensibilidad a la estructura de dependencia en colas bajo esta covarianza nueva: son extensiones pendientes, no efectivamente corridas, y no se reportan cifras sin haberlas corrido.

\section{Limitaciones del Análisis de Proyecciones}
\label{sec:limitaciones_dsa}

Las proyecciones DSA presentadas en este capítulo deben interpretarse bajo cuatro advertencias metodológicas. Primera, la matriz de covarianzas del componente estocástico combina residuos estimados del SVAR con calibraciones históricas para las tasas de interés doméstica y externa, que no tienen serie propia en el sistema de cointegración (Sección \ref{sec:dsa_metodologia} del Capítulo \ref{ch:metodologia}). Segunda, la identidad contable adopta una regla de reacción fiscal pasiva en concordancia con los hallazgos del Capítulo \ref{ch:resultados}. Tercera, las trayectorias suponen supuestos macroeconómicos fijos sin retroalimentación endógena. Cuarta, la distribución $t$ multivariada introduce dependencia en colas (\textit{tail dependence}) que podría sobreestimar la frecuencia de eventos extremos simultáneos; un análisis de sensibilidad a esta propiedad ($\lambda\approx0.076$), corrido sobre la especificación anterior calibrada a mano, desplazó la probabilidad central en un rango moderado de 2 a 7 puntos porcentuales sin revertir ninguna conclusión cualitativa; no se cuenta todavía con el análogo para la especificación de referencia actual.

La sostenibilidad de la deuda pública argentina hacia 2035 depende, bajo ambas especificaciones del componente estocástico ($26{,}1\%$ de referencia, $29{,}8$--$32{,}5\%$ en la especificación anterior calibrada a mano), de la evolución de las variables macroeconómicas fundamentales en un grado mucho mayor que de cualquier ajuste fiscal discrecional marginal, sobre las cuales la política fiscal \textit{per se} ejerce un control solo parcial.
